\section{Estados financieros}

Los estados financieros utilizados para el analisis son:

\begin{itemize}
  \item Estado de Situacion Financiera.
  \item Estado de Resultados.
  \item Estado de Flujos de Efectivo.
  \item Estado de Cambios en el Patrimonio.
\end{itemize}

Para el caso practico se utilizan principalmente el Estado de Situacion Financiera y el Estado de Resultados.

\subsection{Estado de Situacion Financiera}

El Estado de Situacion Financiera muestra los activos, pasivos y patrimonio de la empresa en una fecha determinada. La ecuacion basica es:

\[
\text{Activos} = \text{Pasivos} + \text{Patrimonio}
\]

En el analisis vertical, el total activo suele tomarse como base para evaluar la estructura financiera.

\subsection{Estado de Resultados}

El Estado de Resultados muestra ingresos, costos, gastos, impuesto y utilidad neta durante un periodo. Permite evaluar el desempeno economico y la eficiencia de la empresa para convertir ventas en utilidad.

En el analisis vertical, las ventas netas suelen tomarse como base del 100\%.
